\subsection{Slaughter Trace}\label{sec:trace-slaughter}

We now consider a trace obtained by evaluating the 
displacement and rotation at a fixed reference point $\CenterTrace$ on 
the cross-section (e.g.\ the centroid or shear center). The idea is to use the 
local rigid approximation at $\CenterTrace$ rather than an $L^2$-optimal 
one.

\begin{definition}[Pointwise trace]
\label{def:pointwise-trace}
Let $\CenterTrace \in \Section$ be a fixed reference point. The 
\emph{pointwise trace} associated with $\CenterTrace$ is defined by
\begin{equation}
\label{eq:def-point-trace}
\bm{u}(x) := \hat{\bm{u}}(x,\CenterTrace),
\qquad
\bm{\theta}(x) := \tfrac{1}{2}\,\nabla_{\bm{x}} \times 
\hat{\bm{u}}(x,\bm{r})\big|_{\bm{r}=\CenterTrace}.
\end{equation}
The warping amplitude $\bm{\alpha}(x)$ is determined by the residual
\begin{equation}
\Phi(\bm{r})\,\bm{\alpha}(x)
=
\hat{\bm{u}}(x,\bm{r}) - \bm{u}(x) - \bm{\theta}(x)\times\bm{r}.
\end{equation}
\end{definition}

\begin{remark}
Unlike the section-average trace of \cref{sec:trace}, the pointwise trace does 
not minimize the $L^2$-norm of the warping residual; instead, it matches the 
exact displacement and local rotation at a single reference point 
$\CenterTrace$. The residual 
$\tilde{\bm{u}} := \hat{\bm{u}} - (\bm{u} + \bm{\theta}\times\bm{r})$ 
is generally \emph{not} orthogonal to the section-rigid modes.
\end{remark}

We derive the dependence of $\bm{e}(x)$ on $\bm{p}$ by 
considering each mode separately.

\paragraph{Extension.}

The extension mode displacement is
\begin{equation*}
\hat{\bm{u}}^{(\varepsilon)} = (x\,\FrameBasis - \nu\,\bm{r})\,a_\varepsilon.
\end{equation*}
Evaluating at the reference point $\CenterTrace$:
\begin{equation*}
\bm{u}^{(\varepsilon)}(x)
=
\hat{\bm{u}}^{(\varepsilon)}(x,\CenterTrace)
=
\big(x\,\FrameBasis - \nu\,\CenterTrace\big)\,a_\varepsilon.
\end{equation*}

The rotation field associated with $\hat{\bm{u}}^{(\varepsilon)}$ is obtained 
from the curl. The displacement gradient is
\[
\nabla_{\bm{x}}\hat{\bm{u}}^{(\varepsilon)}
=
\FrameBasis \otimes \FrameBasis - \nu\,\mathbf{1}_\Omega,
\]
whose skew part vanishes. Thus the local rotation vector is identically zero:
\begin{equation*}
\bm{\theta}^{(\varepsilon)}(x) = \mathbf{0}.
\end{equation*}

The strains are therefore
\begin{equation*}
\bm{\gamma}^{(\varepsilon)} 
= (\bm{u}^{(\varepsilon)})' + \FrameBasis \times \bm{\theta}^{(\varepsilon)}
= \FrameBasis\,a_\varepsilon,
\qquad
\bm{\kappa}^{(\varepsilon)} 
= (\bm{\theta}^{(\varepsilon)})' = \mathbf{0}.
\end{equation*}

Hence
\begin{equation}
\bm{e}^{(\varepsilon)}(x)
=
\begin{pmatrix}
\FrameBasis \\[2pt]
\mathbf{0}
\end{pmatrix} a_\varepsilon,
\end{equation}
which is constant in $x$. In particular, the contribution of $a_\varepsilon$ to 
$\bm{e}_0$ is $\FrameBasis$ in the axial strain component and zero elsewhere, 
and the contribution to $\bm{e}_1$ is zero.

\paragraph{Bending.}

The bending mode displacement is
\begin{equation*}
\hat{\bm{u}}^{(a)} = \Big( -\tfrac{1}{2}x^2\,\mathbf{1}_\Omega 
- \nu\,\mathrm{dev}(\bm{r} \otimes \bm{r}) 
+ x\,\FrameBasis \otimes \bm{r} \Big) \bm{a}_\Omega.
\end{equation*}

\subparagraph{Pointwise deflection.}
Evaluating at $\CenterTrace$:
\begin{align*}
\bm{u}^{(a)}(x)
&=
\hat{\bm{u}}^{(a)}(x,\CenterTrace) \\
&=
\Big( -\tfrac{1}{2}x^2\,\mathbf{1}_\Omega 
- \nu\,\mathrm{dev}(\CenterTrace \otimes \CenterTrace) 
+ x\,\FrameBasis \otimes \CenterTrace \Big) \bm{a}_\Omega \\
&=
-\tfrac{1}{2}x^2\,\bm{a}_\Omega
- \nu\,\mathrm{dev}(\CenterTrace \otimes \CenterTrace)\,\bm{a}_\Omega
+ x\,(\CenterTrace\cdot\bm{a}_\Omega)\,\FrameBasis.
\end{align*}

\subparagraph{Pointwise rotation.}
The rotation vector is given by 
$\bm{\theta}^{(a)} = \tfrac{1}{2}\nabla_{\bm{x}}\times\hat{\bm{u}}^{(a)}$ 
evaluated at $\CenterTrace$. The only part contributing 
a nonzero curl is the sum of the $x\,\FrameBasis\otimes\bm{r}$ term and the 
Poisson-bending term. As in the section-average case, this can be written as
\begin{equation}
\bm{\theta}^{(a)}(x)
=
-x\,(\FrameBasis \times \bm{a}_\Omega)
+ \bm{\Xi}^{(a)}_{\theta,\mathrm{pt}}\,\bm{a}_\Omega,
\label{eq:theta-a-pt}
\end{equation}
where the \emph{pointwise Poisson-bending rotation tensor} is defined by
\begin{equation}
\bm{\Xi}^{(a)}_{\theta,\mathrm{pt}} 
:= \tfrac{1}{2}\,\nabla_{\bm{x}}\times\!\big[-\nu\,\mathrm{dev}(\bm{r} \otimes \bm{r})\big]
\big|_{\bm{r}=\CenterTrace}
\in \mathbb{R}^{3\times 2}.
\end{equation}

\subparagraph{Bending strains.}
Differentiating:
\begin{equation*}
\bm{\kappa}^{(a)}(x)
= (\bm{\theta}^{(a)})' 
= -\FrameBasis \times \bm{a}_\Omega,
\end{equation*}
which is constant in $x$.

For the shear-type strain,
\begin{align*}
(\bm{u}^{(a)})'(x)
&=
-x\,\bm{a}_\Omega + (\CenterTrace\cdot\bm{a}_\Omega)\,\FrameBasis, \\
\bm{\gamma}^{(a)}(x)
&=
(\bm{u}^{(a)})' + \FrameBasis \times \bm{\theta}^{(a)} \\
&=
-x\,\bm{a}_\Omega + (\CenterTrace\cdot\bm{a}_\Omega)\,\FrameBasis 
+ \FrameBasis \times\big(-x\,(\FrameBasis \times \bm{a}_\Omega) 
+ \bm{\Xi}^{(a)}_{\theta,\mathrm{pt}}\,\bm{a}_\Omega\big) \\
&=
-x\,\bm{a}_\Omega + (\CenterTrace\cdot\bm{a}_\Omega)\,\FrameBasis 
+ x\,\bm{a}_\Omega + \FrameBasis \times \bm{\Xi}^{(a)}_{\theta,\mathrm{pt}}\,\bm{a}_\Omega \\
&=
(\CenterTrace\cdot\bm{a}_\Omega)\,\FrameBasis 
+ (\FrameBasis \times \bm{\Xi}^{(a)}_{\theta,\mathrm{pt}})\,\bm{a}_\Omega.
\end{align*}
Using $\bm{a}_\Omega \perp \FrameBasis$, the $x$-dependent terms cancel as in 
the section-average trace. The remaining constant term can be grouped into a 
tensor acting on $\bm{a}_\Omega$. In particular, if 
$\CenterTrace$ lies on the centroidal axis (so that 
$\CenterTrace\cdot\bm{a}_\Omega = 0$), then
\begin{equation*}
\bm{\gamma}^{(a)}(x)
=
(\FrameBasis \times \bm{\Xi}^{(a)}_{\theta,\mathrm{pt}})\,\bm{a}_\Omega,
\end{equation*}
which is constant in $x$.

Thus, in general,
\begin{equation}
\bm{e}^{(a)}(x)
=
\begin{pmatrix}
\bm{\gamma}^{(a)} \\
\bm{\kappa}^{(a)}
\end{pmatrix}
=
\begin{pmatrix}
\bm{\Gamma}^{(a)} \\
-\FrameBasis \times \mathbf{1}_\Omega
\end{pmatrix}\bm{a}_\Omega,
\end{equation}
with 
$\bm{\Gamma}^{(a)} := 
(\CenterTrace\cdot \cdot)\,\FrameBasis 
+ \FrameBasis \times \bm{\Xi}^{(a)}_{\theta,\mathrm{pt}}$ an appropriate 
$3\times 2$ tensor, and $\bm{e}^{(a)}(x)$ is constant in $x$.

\paragraph{Torsion.}

The torsion mode displacement is
\begin{equation*}
\hat{\bm{u}}^{(\vartheta)} = \big( x\,(\FrameBasis \times \bm{r}) 
+ \WarpTwistIesan(\bm{r})\,\FrameBasis \big)\,a_\vartheta.
\end{equation*}

\subparagraph{Pointwise deflection.}
Evaluating at $\CenterTrace$:
\begin{equation*}
\bm{u}^{(\vartheta)}(x)
=
\big( x\,(\FrameBasis \times \CenterTrace) 
+ \WarpTwistIesan(\CenterTrace)\,\FrameBasis \big)\,a_\vartheta.
\end{equation*}

\subparagraph{Pointwise rotation.}
The rigid part $x\,(\FrameBasis \times \bm{r})$ corresponds to a pure twist 
about $\FrameBasis$ with rotation vector $x\,\FrameBasis$; the warping term 
$\WarpTwistIesan\,\FrameBasis$ contributes an additional $x$-independent 
rotation at $\CenterTrace$. Thus
\begin{equation}
\bm{\theta}^{(\vartheta)}(x)
=
\big(x\,\FrameBasis + \CenterTwistLocation\big)\,a_\vartheta,
\label{eq:theta-torsion-pt}
\end{equation}
where we define the \emph{pointwise twist center}
\begin{equation}
\CenterTwistLocation
:=
\tfrac{1}{2}\,\nabla_{\bm{x}}\times\!\big(\WarpTwistIesan(\bm{r})\,\FrameBasis\big)
\big|_{\bm{r}=\CenterTrace}
\in \mathbb{R}^3.
\end{equation}

\subparagraph{Torsional strains.}
Differentiating:
\begin{equation*}
\bm{\kappa}^{(\vartheta)}(x)
=
(\bm{\theta}^{(\vartheta)})' = \FrameBasis\,a_\vartheta,
\end{equation*}
constant in $x$.

The shear-type strain is
\begin{align*}
(\bm{u}^{(\vartheta)})'(x)
&=
(\FrameBasis \times \CenterTrace)\,a_\vartheta, \\
\bm{\gamma}^{(\vartheta)}(x)
&=
(\bm{u}^{(\vartheta)})' + \FrameBasis \times \bm{\theta}^{(\vartheta)} \\
&=
(\FrameBasis \times \CenterTrace)\,a_\vartheta
+ \FrameBasis \times (x\,\FrameBasis + \CenterTwistLocation)\,a_\vartheta \\
&=
(\FrameBasis \times \CenterTrace)\,a_\vartheta
+ \FrameBasis \times \CenterTwistLocation\,a_\vartheta,
\end{align*}
since $\FrameBasis \times \FrameBasis = \mathbf{0}$. 
%
Thus
\begin{equation}
\bm{e}^{(\vartheta)}(x)
=
\begin{pmatrix}
\FrameBasis \times \CenterTrace 
+ \FrameBasis \times \CenterTwistLocation \\
\FrameBasis
\end{pmatrix} a_\vartheta,
\end{equation}
which is constant in $x$. In particular, for $\CenterTrace = \mathbf{0}$ 
(centroidal point), the shear part reduces to 
$\FrameBasis \times \CenterTwistLocation\,a_\vartheta$.

\paragraph{Flexure.}

The flexure mode displacement is
\begin{equation*}
\begin{aligned}
\hat{\bm{u}}^{(b)} &= 
\big(x\,\FrameBasis \times \bm{r} + \WarpTwistIesan(\bm{r})\,\FrameBasis\big)\,c_\vartheta 
\\
&\quad
+ \Big( -\tfrac{1}{6}x^3\,\mathbf{1}_\Omega 
+ x\,\bm{\Pi}_{(b)}(\bm{r}) 
+ \tfrac{1}{2}x^2\,\FrameBasis \otimes (\bm{r} - \CentroidLocation) 
+ \FrameBasis \otimes \bm{\WarpShearIesan}(\bm{r}) \Big)\,\bm{b}_{\Section},
\end{aligned}
\end{equation*}
where $c_\vartheta = \CenterShearLocation \cdot \bm{b}_{\Section}$ is the 
flexure-induced twist with shear center $\CenterShearLocation$ as in 
\eqref{eq:def-shear-center}.

\subparagraph{Pointwise deflection.}
Evaluating at $\CenterTrace$:
\begin{align*}
\bm{u}^{(b)}(x)
&=
\big(x\,\FrameBasis \times \CenterTrace 
+ \WarpTwistIesan(\CenterTrace)\,\FrameBasis\big)\,c_\vartheta \\
&\quad
+ \Big( -\tfrac{1}{6}x^3\,\mathbf{1}_\Omega 
+ x\,\bm{\Pi}_{(b)}(\CenterTrace) 
+ \tfrac{1}{2}x^2\,\FrameBasis \otimes (\CenterTrace - \CentroidLocation) 
+ \FrameBasis \otimes \bm{\WarpShearIesan}(\CenterTrace) \Big)\,\bm{b}_{\Section}.
\end{align*}
Expanding:
\begin{align}
\bm{u}^{(b)}(x)
&=
-\tfrac{1}{6}x^3\,\bm{b}_{\Section}
+ x\,\bm{\Pi}_{(b)}(\CenterTrace)\,\bm{b}_{\Section}
+ \tfrac{1}{2}x^2\,\FrameBasis\,\big((\CenterTrace - \CentroidLocation)\cdot\bm{b}_{\Section}\big)
+ \FrameBasis\,\big(\bm{\WarpShearIesan}(\CenterTrace)\cdot\bm{b}_{\Section}\big)
\notag\\
&\quad
+ \big(x\,\FrameBasis \times \CenterTrace 
+ \WarpTwistIesan(\CenterTrace)\,\FrameBasis\big)\,(\CenterShearLocation \cdot \bm{b}_{\Section}).
\label{eq:u-b-pt}
\end{align}

\subparagraph{Pointwise rotation.}
Similarly, the rotation vector at $\CenterTrace$ can be written in the 
general form (by polynomial structure and symmetry)
\begin{equation}
\bm{\theta}^{(b)}(x)
=
-\tfrac{1}{2}x^2\,(\FrameBasis \times \bm{b}_{\Section})
+ x\,\bm{A}_{\theta,1}^{\mathrm{pt}}\,\bm{b}_{\Section}
+ \bm{A}_{\theta,0}^{\mathrm{pt}}\,\bm{b}_{\Section},
\label{eq:theta-b-pt}
\end{equation}
where the \emph{pointwise flexural rotation tensors} 
$\bm{A}_{\theta,1}^{\mathrm{pt}},\bm{A}_{\theta,0}^{\mathrm{pt}}\in\mathbb{R}^{3\times 2}$ 
collect contributions from the flexure-induced twist $c_\vartheta$, the shear 
warping $\bm{\Pi}_{(b)}$, and the Ieşan shear warping $\bm{\WarpShearIesan}$ 
evaluated at $\CenterTrace$. Their explicit form is not needed for the 
polynomial structure of the strains.

Differentiating:
\begin{equation}
\bm{\kappa}^{(b)}(x)
=
(\bm{\theta}^{(b)})' 
=
-x\,(\FrameBasis \times \bm{b}_{\Section})
+ \bm{A}_{\theta,1}^{\mathrm{pt}}\,\bm{b}_{\Section},
\label{eq:kappa-b-pt}
\end{equation}
which is affine in $x$.

\subparagraph{Flexural shear strain.}
Differentiating \eqref{eq:u-b-pt}:
\begin{align*}
(\bm{u}^{(b)})'(x)
&=
-\tfrac{1}{2}x^2\,\bm{b}_{\Section}
+ \bm{\Pi}_{(b)}(\CenterTrace)\,\bm{b}_{\Section}
+ x\,\FrameBasis\,\big((\CenterTrace - \CentroidLocation)\cdot\bm{b}_{\Section}\big) \\
&\quad
+ \big(\FrameBasis \times \CenterTrace\big)\,(\CenterShearLocation \cdot \bm{b}_{\Section}).
\end{align*}
Thus
\begin{align*}
\bm{\gamma}^{(b)}(x)
&=
(\bm{u}^{(b)})' + \FrameBasis \times \bm{\theta}^{(b)} \\
&=
\Big(-\tfrac{1}{2}x^2\,\bm{b}_{\Section} 
+ \tfrac{1}{2}x^2\,\bm{b}_{\Section}\Big) \\
&\quad
+ x\Big(\FrameBasis \otimes (\CenterTrace - \CentroidLocation) 
+ \FrameBasis \times \bm{A}_{\theta,1}^{\mathrm{pt}}\Big)\,\bm{b}_{\Section} \\
&\quad
+ \Big(\bm{\Pi}_{(b)}(\CenterTrace) 
+ (\FrameBasis \times \CenterTrace)\otimes \CenterShearLocation 
+ \FrameBasis \times \bm{A}_{\theta,0}^{\mathrm{pt}}\Big)\,\bm{b}_{\Section}.
\end{align*}
The $x^2$-terms cancel, as expected. Define the pointwise flexural shear tensors
\begin{align}
\bm{\Gamma}^{\mathrm{pt}}_1 
&:= \FrameBasis \otimes (\CenterTrace - \CentroidLocation) 
+ \FrameBasis \times \bm{A}_{\theta,1}^{\mathrm{pt}},
\label{eq:def-Gamma1-pt} \\
\bm{\Gamma}^{\mathrm{pt}}_0 
&:= \bm{\Pi}_{(b)}(\CenterTrace) 
+ (\FrameBasis \times \CenterTrace)\otimes \CenterShearLocation 
+ \FrameBasis \times \bm{A}_{\theta,0}^{\mathrm{pt}}.
\label{eq:def-Gamma0-pt}
\end{align}

Then
\begin{equation}
\bm{\gamma}^{(b)}(x)
=
x\,\bm{\Gamma}^{\mathrm{pt}}_1\,\bm{b}_{\Section}
+ \bm{\Gamma}^{\mathrm{pt}}_0\,\bm{b}_{\Section},
\label{eq:gamma-flexure-pt}
\end{equation}
which is affine in $x$, and together with \eqref{eq:kappa-b-pt} yields the 
flexural contribution to $\bm{e}(x)$.

\paragraph{Assembly of trace matrices.}

Collecting the contributions from all four Ieşan modes, the pointwise beam 
strains can be written in the affine form
\begin{equation}
\bm{e}(x)
=
\bm{e}_{0,\mathrm{pt}} + x\,\bm{e}_{1,\mathrm{pt}}
=
\mathbf{T}_{0,\mathrm{pt}}\,\bm{p}
+ x\,\mathbf{T}_{1,\mathrm{pt}}\,\bm{p},
\end{equation}
where $\bm{p} = (a_\varepsilon, \bm{a}_\Omega, a_\vartheta, \bm{b}_{\Section})^\top$ 
and $\bm{e} = (\epsilon, \bm{\gamma}_\perp, \varkappa, \bm{\kappa}_\perp)^\top$.

From the mode-by-mode calculations above, the \emph{pointwise trace matrices} 
$\mathbf{T}_{0,\mathrm{pt}},\mathbf{T}_{1,\mathrm{pt}} \in \mathbb{R}^{6\times 6}$ 
have the block structure
\begin{equation}
\mathbf{T}_{0,\mathrm{pt}} = 
\begin{pmatrix}
1 & (\bm{\Gamma}^{(a)})_\parallel^\top & 0 & (\bm{\Gamma}^{\mathrm{pt}}_0)_\parallel^\top \\[4pt]
\mathbf{0} & (\bm{\Gamma}^{(a)})_\perp & (\FrameBasis \times \CenterTwistLocation + \FrameBasis \times \CenterTrace)_\perp & (\bm{\Gamma}^{\mathrm{pt}}_0)_\perp \\[4pt]
0 & 0 & 1 & (\bm{A}_{\theta,1}^{\mathrm{pt}})_\parallel^\top \\[4pt]
\mathbf{0} & -\FrameBasis^{\times} & \mathbf{0} & (\bm{A}_{\theta,1}^{\mathrm{pt}})_\perp
\end{pmatrix},
\label{eq:T0-pointwise}
\end{equation}
\begin{equation}
\mathbf{T}_{1,\mathrm{pt}} = 
\begin{pmatrix}
0 & \mathbf{0}^\top & 0 & (\bm{\Gamma}^{\mathrm{pt}}_1)_\parallel^\top \\[4pt]
\mathbf{0} & \mathbf{0} & \mathbf{0} & (\bm{\Gamma}^{\mathrm{pt}}_1)_\perp \\[4pt]
0 & 0 & 0 & (-\FrameBasis \times \mathbf{1}_\Omega)_\parallel^\top \\[4pt]
\mathbf{0} & \mathbf{0} & \mathbf{0} & (-\FrameBasis \times \mathbf{1}_\Omega)_\perp
\end{pmatrix},
\label{eq:T1-pointwise}
\end{equation}
where $(\cdot)_\parallel$ and $(\cdot)_\perp$ denote the axial and transverse 
components with respect to $\FrameBasis$, and the various tensors
$\bm{\Xi}^{(a)}_{\theta,\mathrm{pt}}$, 
$\CenterTwistLocation$, 
$\bm{A}_{\theta,0}^{\mathrm{pt}}$, 
$\bm{A}_{\theta,1}^{\mathrm{pt}}$, 
$\bm{\Gamma}^{(a)}$, 
$\bm{\Gamma}^{\mathrm{pt}}_0$, 
$\bm{\Gamma}^{\mathrm{pt}}_1$
are as defined above.

\begin{remark}[Warping index of the pointwise trace]
For a section with non-degenerate geometry (invertible inertia tensor and 
nontrivial warping fields), the six columns of the stacked matrix
\begin{equation}
\mathbf{T}_{\mathrm{stack,pt}}
:=
\begin{pmatrix}
\mathbf{T}_{0,\mathrm{pt}} \\[2pt]
\mathbf{T}_{1,\mathrm{pt}}
\end{pmatrix}
\end{equation}
are generically linearly independent, yielding 
$\mathrm{rank}(\mathbf{T}_{\mathrm{stack,pt}})=6$. 
In this case, the full Saint-Venant parameter vector 
$\bm{p}$ is uniquely determined by the pointwise beam strains 
$\bm{e}(x)$ (equivalently, by their coefficients 
$\bm{e}_{0,\mathrm{pt}},\bm{e}_{1,\mathrm{pt}}$), and a six-strain Timoshenko 
beam with this trace can, in principle, reproduce the entire classical 
Saint-Venant torsion–extension–shear–bending class.
\end{remark}
